\title{LongRoPE: Extending LLM Context Window Beyond 2 Million Tokens}

\begin{abstract}
Expanding the context window is an important capability for large language models (LLMs). Nevertheless, current models with longer context windows are typically restricted to approximately 128k tokens, largely because of expensive fine-tuning requirements, limited availability of lengthy textual data, and disruptive effects arising from introducing new token positions.

In this work, we present LongRoPE, a novel approach that, for the first time, pushes the context window of existing pre-trained LLMs to a remarkable \textbf{2048k} tokens, requiring as few as 1k fine-tuning steps limited to training sequences of up to 256k tokens, all while preserving the model’s effectiveness on short contexts. This advance is enabled by three main contributions: (i) we systematically discover and utilize two distinct non-uniformities inherent in positional interpolation via an efficient search process, which allows for superior initialization during fine-tuning and supports up to an 8$\times$ extension of the context window even without fine-tuning; (ii) we propose a progressive expansion method where the model is first fine-tuned at the 256k token length, followed by a subsequent positional interpolation on this longer-context LLM, achieving the 2048k window; (iii) we recalibrate LongRoPE at the 8k token length to restore and maintain the performance seen at the original short context range. Comprehensive empirical evaluation on both LLaMA2 and Mistral, spanning diverse benchmark tasks, confirms the advantages of our approach. LongRoPE-augmented models maintain the core architecture, modifying only the positional embeddings with minimal changes, allowing them to retain compatibility with most existing optimizations. The code will be made accessible at \texttt{https://github.com/microsoft/LongRoPE}

\section{Introduction}

Although Large Language Models (LLMs) have demonstrated impressive capabilities across a range of tasks~\cite{gpt4,llama2}, they are constrained by restricted context window sizes, such as the 4096-token limit of LLaMA2~\cite{llama2}. When input exceeds this window, the performance of LLMs deteriorates, primarily because the models have not been trained to handle these additional sequence positions. This limitation creates difficulties in key applications, including in-context learning with multiple examples~\cite{Huang2023FewerIM} and the operation of LLM agents~\cite{park2023generative,madaan2023selfrefine}.

Recent studies have demonstrated that the context window of pre-trained LLMs can be pushed up to approximately 128k tokens by employing fine-tuning techniques with lengthier documents~\cite{longlora,pi,yarn,zhang2024soaring,liu2023scaling}. However, scaling the context window further faces three principal hurdles. Firstly, introducing new positional indices that have not been trained results in numerous anomalous values, exacerbating out-of-distribution phenomena and hindering the convergence of the fine-tuning process~\cite{pi}. This becomes significantly problematic when expanding from 4k to beyond 1000k, where over 90\% of positions are novel. Secondly, fine-tuning generally relies on access to texts matching these extended lengths, yet corpora containing such ultra-long sequences, especially those surpassing 1000k, are scarce. Additionally, training on these exceptionally long texts imposes a substantial computational burden, demanding substantial time and GPU capacity. Thirdly, with extremely extended context windows, the attention mechanism disperses across the vastly increased token space, thereby diminishing the model's efficacy on tasks requiring short-range context~\cite{pi}.

To address the initial challenge, a common strategy involves interpolating RoPE positional embedding~\cite{rope,pi}, which adjusts new position indices so that they fit within the scope of the original pretrained range, illustrated in Fig.\ref{fig:overview}. Position Interpolation (PI)~\cite{pi} modifies the rotary angles in RoPE through linear interpolation, scaling them by the extension ratio. NTK~\cite{ntk,dynamicntk} recommends varying interpolation and extrapolation across the RoPE's dimensions rather than treating them equally. YaRN~\cite{yarn} takes a frequency-based approach by grouping RoPE dimensions into three categories and then applying extrapolation, NTK-style, and linear interpolation correspondingly. Nonetheless, the information entropy associated with positional embedding within the Transformer framework is inherently complex and non-uniform. This nuanced heterogeneity is largely overlooked in current methods, resulting in information degradation and restricting the achievable context window.

Section~\ref{sec:analysis} empirically demonstrates two principal insights: \textbf{(1)} High-quality positional interpolation must address two distinct non-uniformities—variations across RoPE dimensions as well as across token positions. Dimensions associated with lower RoPE values and earlier token positions are less sensitive to interpolation; nevertheless, the ideal interpolation patterns shift with different target sequence lengths. 
\textbf{(2)} Incorporating these identified non-uniformities during positional interpolation enables more faithful preservation of information present in the original RoPE embeddings, especially regarding critical dimensions and positions. This strategy mitigates information loss due to positional interpolation, thereby enhancing the initialization quality for subsequent fine-tuning. Additionally, it enables extending the context length by a factor of 8$\times$ even in the absence of fine-tuning.

Building upon these results, we present \textbf{LongRoPE}, a robust approach designed to enlarge the LLM context window to accommodate over 2 \emph{million} tokens. The LongRoPE methodology is underpinned by three principal innovations. To start, {LongRoPE} leverages the multidimensional aspects of non-uniform positional interpolation. Specifically, it determines optimal rescale factors for the rotation angles in RoPE corresponding to each dimension, tailored to the token positions. Since the complexity of locating appropriate rescale factors grows exponentially as the desired extension ratio increases, {LongRoPE} incorporates an evolutionary search strategy enhanced by two specialized optimization techniques to improve search performance. An illustration of the resulting rescaled RoPE found through this process is depicted in Fig.~\ref{fig:overview}.

Subsequently, {LongRoPE} utilizes an effective incremental extension approach, enabling a context window of up to 2048k without requiring direct fine-tuning on ultra-long documents, which are both rare and challenging to obtain. The process initiates by identifying a 256k token length using the pre-trained LLM, followed by fine-tuning at this intermediate length. Because our non-uniform positional interpolation mechanism permits up to an 8$\times$ increase in context without additional fine-tuning, we proceed with a secondary search to identify new RoPE rescaling parameters on the fine-tuned, extended LLM. Through this methodology, a 2048k context window is successfully attained for both LLaMA2 and Mistral~\cite{mistral}.

Lastly, to prevent a decline in performance when using the initial (smaller) context window, {LongRoPE} further tunes the RoPE rescale parameters in the expanded LLM. Using the same search method applied for scaling up from 256k to 2048k, we also scale down to 4k and 8k context windows on the LLM fine-tuned at 256k, aiming to reduce positional interpolation. For inference, when the input sequence is under 8k tokens, RoPE is modified with the optimized rescale factors identified by our search algorithm.

A comprehensive set of experiments conducted on multiple LLMs and a range of long-context tasks validates the efficacy of our proposed approach. Our findings indicate that LongRoPE reliably sustains low perplexity across evaluation lengths spanning from 4k up to 2048k, attains passkey retrieval accuracy exceeding 90\%, and matches the performance of baseline models on established benchmarks with a 4096 token context window. The LongRoPE technique is broadly compatible with any LLM architecture utilizing RoPE embeddings. We intend to make both our implementation and the LongRoPE-2048k models publicly available.

\section{Non-uniformity in Positional Interpolation}
\label{sec:analysis}

\subsection{Preliminary}

Transformer models require explicit positional information, often in the form of position embedding, to represent the order of input tokens. Our work focuses on the RoPE~\cite{rope} position embedding, which is widely used in recent  LLMs. For a token at position index $n$, its corresponding RoPE encoding can be simplified as follows:

\begin{equation}
		\fontsize{7.8}{7.8} \selectfont
	\label{eq:rope}
	\left[ cos(n\theta_0), sin(n\theta_0), cos(n\theta_1), \cdot\cdot\cdot, cos(n\theta_{d/2-1}), sin(n\theta_{d/2-1}) \right]   
\end{equation}

where $d$ stands for the embedding dimensionality, $n\theta_i$ denotes the rotary phase assigned to the token in position $n$, and $\theta_i=\theta^{-2i/d}$ captures the frequency of rotation. Within RoPE, $\theta$ is set to a default value of 10000.

\textbf{\emph{Context window extension ratio $s$} and positional interpolation}. The parameter $s$ is defined as the proportion of the new, expanded context length $L'$ relative to the initial length $L$, given by $s=\frac{L'}{L}$.

To extend the context window from $L$ to $L'$, current positional interpolation methods suggest downscaling rotation frequencies $\theta_i$ based on extension ratio $s$. Let $\beta=\theta^{2/d}$, and $\lambda$ denote the actual rescale factor related to $s$, we   unify these positional interpolation methods as follows:

\begin{equation}	
	\fontsize{7.5}{7.5} \selectfont
	\label{eq:unified_rope}
	\begin{aligned}
		\left[cos\left(\frac{n}{\lambda(\beta)^0}\right), sin \left(\frac{n}{\lambda(\beta)^0}\right), cos\left(\frac{n}{\lambda(\beta)^1}\right), \cdot\cdot\cdot,  sin \left(\frac{n}{\lambda(\beta)^{d/2-1}}\right) \right]   
	\end{aligned}
\end{equation}

\textbf{Linear positional interpolation (PI)}. PI~\cite{pi} introduces a method where position indices are interpolated linearly as long as they remain within the original training sequence length. To expand the context window by a factor of $s$, each positional rotation angle is uniformly scaled down by $\lambda = s$ across every RoPE dimension. While this approach allows for longer input sequences, it results in compressed or "crowded" positional encodings, making it challenging for the model to differentiate between tokens that are near each other in position. Consequently, PI typically exhibits lower performance when applied with large extension ratios.

\textbf{NTK-based interpolation and extrapolation}. ~\cite{ntk,dynamicntk} approach RoPE from the standpoint of information encoding, leveraging Neural Tangent Kernel (NTK) theory~\cite{jacot2018neural,tancik2020fourier} in their analysis. To address the challenge of crowded positions inherent to PI, they advocate redistributing interpolation across RoPE dimensions. This method moderates the scaling for lower (high frequency) dimensions while increasing it for higher (low frequency) ones, thus enabling both interpolation and extrapolation for positional encodings, formalized as $\lambda=s^{i}$. The advanced dynamic NTK method~\cite{dynamicntk} further modifies the extension ratio on a position-by-position basis, reflecting the current sequence length. Distinct from PI—which requires fine-tuning—NTK-based techniques are capable of extending the context window even without fine-tuning, typically supporting an extension ratio up to 4$\times$.

\textbf{YaRN}~\cite{yarn} represents a notable advance in the effectiveness of positional interpolation. This method partitions the RoPE dimensions according to their frequency characteristics into three separate categories, each assigned a distinct interpolation approach. Dimensions with high frequency are subject to extrapolation ($\lambda$=1), while those at the low-frequency end utilize linear interpolation (PI). For the remaining, intermediate-frequency RoPE dimensions, the NTK method is employed. Central to YaRN’s design is this dimension grouping mechanism, which is determined through manual, empirically-driven experimentation. Consequently, this reliance on human decision-making could hinder optimal results when applied to novel LLM architectures.

\subsection{Study on Non-uniform Positional Interpolation}

Drawing motivation from NTK and YaRN, we observe that their improvements stem from leveraging non-linearities, particularly by evaluating varying frequencies along RoPE dimensions to achieve tailored interpolation and extrapolation. Despite these advances, existing non-linear mechanisms are predominantly governed by manually crafted guidelines. This observation brings forth two pertinent questions: (1) Does the present method for positional interpolation represent the best possible approach? (2) Might there be other non-linearities yet to be investigated?

In order to investigate these issues, we employ evolution search (refer to Sec.~\ref{sec:method}) to identify improved non-uniform positional interpolation schemes for LLaMA2-7B. The search process is steered by evaluating perplexity on 5 randomly selected examples from the PG19~\cite{pg19} validation dataset. Our experiments yield several important insights, as outlined below.

\textbf{Finding 1}: \textit{RoPE dimensions exhibit substantial non-uniformities, which are not effectively handled by current positional interpolation methods.}

We determine the optimal value of $\lambda$ for each RoPE dimension as presented in Eq.~\ref{eq:unified_rope}. Table~\ref{tbl:exp1} provides a comparison of LLaMA2-7B’s perplexity across various approaches on the PG19 and Proof-pile~\cite{proof-pile} test sets, evaluated without additional fine-tuning. Our method, based on tailored $\lambda$ values, demonstrates notable performance gains, highlighting that prevalent linear (PI) and non-uniform techniques (Dynamic-NTK and YaRN) do not deliver optimal results. It is worth emphasizing that, on PG19, YaRN yields worse results compared to PI and NTK, primarily because it fails to cover the full intended context window in the absence of fine-tuning. Specifically, for an 8k context, YaRN's perplexity escalates sharply beyond 7k tokens.

In our investigation, we observed that the adjusted scaling parameters $\lambda$ in Eq.~\ref{eq:unified_rope} exhibit variability, contrasting with the constant scale $s$ utilized in PI, the formulaic approach of NTK, and YaRN's method of applying group-wise scaling. These variable scaling factors yield notable enhancements in LLaMA2's language modeling abilities, specifically in terms of perplexity, when tested on context windows of 8k and 16k tokens without requiring any additional fine-tuning. This improvement arises from the manner in which the modified positional embeddings retain the characteristics of the original RoPE, most importantly in the principal dimensions, thereby alleviating the model’s challenge in identifying tokens with nearby positions.

\textbf{Finding 2}: \textit{RoPE for the initial tokens in the input sequence should be extrapolated with less interpolation.} 

We propose that for the first $\hat{n}$ tokens in each input sequence, the RoPE mechanism should minimize interpolation. The rationale behind this is that these tokens tend to attract substantial attention scores, thereby becoming pivotal within attention layers—a phenomenon reported in Streaming LLM~\cite{streamingllm} and LM-Infinite~\cite{han2023lminfinite}. To examine this assumption, we enlarge the context window to 8k and 16k using PI and NTK, explicitly maintaining the absence of interpolation for the initial $\hat n$ tokens, where $\hat n$ ranges from 0, 2, ..., 256. When $\hat n$=0, the system corresponds to the standard PI and NTK approaches. Table~\ref{tbl:tokenposition} reveals two important findings: \textbf{(1)} withholding interpolation for initial tokens leads to enhanced performance, and \textbf{(2)} the optimal value for $\hat n$ is contingent upon the specific context window extension being targeted.

\textbf{Finding 3}: \textit{Non-uniform positional interpolation effectively extends LLM context window in both fine-tuning and non-fine-tuning settings.} 

Although our experiments have demonstrated that utilizing our searched non-uniform position interpolation substantially enhances extension performance to 8k and 16k contexts in the absence of fine-tuning, achieving effective extension to even larger windows necessitates additional fine-tuning. To this end, we fine-tune LLaMA2-7B incorporating our optimized RoPE on a 64k context window (detailed settings are available in the Appendix). As reported in Table~\ref{tbl:finetune}, our approach delivers marked improvements over both PI and YaRN, regardless of whether LLaMA2-7B is fine-tuned. This performance advantage stems from the strategic application of non-uniform positional interpolation, which minimizes loss of information and yields a superior starting point for subsequent fine-tuning.

\textbf{Summary}. This work identifies two types of non-uniformity—variation across RoPE dimensions and token positions. Harnessing these non-uniformities in positional interpolation leads to substantial advancements in large language model context window extension.

\section{LongRoPE}

\label{sec:method}
Building upon these insights, we propose LongRoPE. This method begins by integrating a streamlined search algorithm designed to leverage both types of non-uniformity effectively, subsequently enabling the expansion of the LLM context window to exceed 2 million tokens.

\subsection{Problem Formulation}

The presence of these two types of non-uniformity significantly enlarges the solution space and adds layers of difficulty to the optimization process. To manage this, we conceptualize the optimization of multidimensional non-uniform position interpolation as a search problem.

For a LLM targeting a  context window size of $L'$ and lengthy input documents $\mathbf{X}$, where each $\mathbf{x}\in\mathbf{X}$ surpasses $L'$ in token length, we denote the original rotary angle of the $i^{th}$ dimension in RoPE embedding at token position $n$ as  $\frac{n}{\beta^i}$. The optimization problem is then formulated as follows:

\begin{equation}
\begin{gathered}
		\label{eq:problem}

\underset {\mathbf{x}\in\mathbf{X}; \, |\mathbf{x}|\ge L'}{ \operatorname {arg\,min} } \,  \mathcal{L}\, \left( \text{LLM} (\text{RoPE}, \mathbf{X})\right), \text{with}  
\\
\scalebox{0.93}{
$\underset {\begin{subarray}{l} i=0,\ldots,\frac{d}{2}-1;\\ n\in[0,|\mathbf{x}|);\end{subarray}}{\text{RoPE}_{(n)}} = \left[\ldots, \cos\left(\mathbb{I}(\hat{\lambda}_{i}, \hat{n}) \cdot \frac{n}{\beta^i}\right), \sin\left(\mathbb{I}(\hat{\lambda}_{i}, \hat{n}) \cdot \frac{n}{\beta^i}\right), \ldots\right]$
}
\\
\text{where } \mathbb{I}(\hat{\lambda}_{i},\hat{n})=
\begin{cases}
1 &\text{if } n < \hat{n} \\
\frac{1}{\lambda_{i}} &\text{if } n \geq \hat{n}
\end{cases}
\end{gathered}
\end{equation}

In this approach, we define a group of rescaling factors, $\mathbb{I}(\hat{\lambda}_{i},\hat{n})$, to accommodate two distinct types of non-uniformities. Here, $\hat{\lambda}_i$ captures the non-uniformity across RoPE dimensions, and $\hat{n}$ represents the threshold for token position non-uniformity. The function $\mathbb{I}(\hat{\lambda}_{i},\hat{n})$ adjusts the rotation angle associated with the $i^{th}$ dimension in RoPE, utilizing $\hat{\lambda}_i$ as the scaling parameter and $\hat{n}$ as the position threshold. For the first $\hat{n}$ minus one token positions, the scaling parameter $\hat{\lambda}_i$ is not active, so the original RoPE rotary angle $\frac{n}{\beta^i}$ is maintained. Conversely, when the token position reaches or exceeds $n \ge \hat{n}$, the rescaling factor is enforced.

With a target context window size of $L'$, our aim is to identify the set of optimal rescale factors—specifically, $\mathbb{I}(\hat{\lambda}_{0},\hat{n})$, $\mathbb{I}(\hat{\lambda}_{1},\hat{n})$, ..., $\mathbb{I}(\hat{\lambda}_{i},\hat{n})$, and so forth—across the $1^{st}$ through $d^{th}$ dimensions of RoPE. The goal is for the resulting $\text{LLM}$, utilizing this rescaled $\text{RoPE}$, to attain the lowest possible next token prediction loss, $\mathcal{L}$ (which is also reflected in its perplexity), when processing input samples $\mathbf{X}$ of token length $L'$.

\subsection{Searching the Non-uniform Position Interpolation}

In addressing the challenge presented in Eq.~\ref{eq:problem}, we propose an efficient approach that identifies the ideal RoPE rescale factors. This technique is specifically designed to maximize the benefits of the multidimensional irregularities within position embedding.

\textbf{Search space}. Our approach involves constructing an extensive search space to encompass the two identified non-uniformities. The structure of this search space is depicted in Table~\ref{tbl:searchspace}. In detail, we permit the determination of individualized rescale factors for each dimension of the RoPE embedding. For the sake of simplifying the search space configuration, we perform the search over $\lambda_i$ and $\hat{n}$ rather than directly optimizing $\mathbb{I}(\hat{\lambda}_{i},\hat{n})$, noting that $\hat{\lambda}_{i}=1/\lambda_i$. According to Table~\ref{tbl:searchspace}, the allowed range for $\lambda_i$ extends from a lower bound of 1.0 (representing direct extrapolation) up to an upper limit of $s\times1.25$ (providing greater interpolation than PI), with increments of 0.01. Here, $s$ denotes the extension ratio for the intended context window.

$\hat{n}$ determines how many leading token positions remain unchanged, foregoing position interpolation in favor of the native RoPE embedding. Based on observed results, $\hat{n}$ is permitted to take values from the set \{0, 1, 2, 4, 8, 12, 16, 20, 24, 28, 32, 64, 128, 256\} during hyperparameter search. If $\hat{n}=0$, every token position applies the respective rescale factors identified by the search.

\textbf{Evolution-based search}. The search space delineated in Table~\ref{tbl:searchspace} encompasses a vast array of positional interpolation configurations, resulting in considerable complexity for systematic exploration. For instance, selecting a $s=4\times$ extension yields $400^{128/2}\times14 = 4\times10^{167}$ possible alternatives. The search space increases at an exponential rate as the extension ratio grows. To effectively navigate this expansive space, we adopt an evolution search strategy~\cite{guo2020single} and integrate two optimization methods designed to markedly improve search efficiency. The full search workflow is depicted in Algorithm~\ref{alg:search}.

\textit{Optimized initial population generation}. In place of purely random initialization for the $P$ rescale factors comprising the population, we ensure inclusion of three specific individuals: the RoPE rescale factors associated with PI, NTK, and YaRN. The rest of the population, consisting of $P$-3 individuals, is created by randomly applying mutations to these three rescale factors with a mutation probability $p$.

\textit{Monotonically non-decreasing constraint}. Once the initial population has been created, LLM perplexity is evaluated for every candidate. This is achieved by assigning the relevant RoPE rescale factors to the target LLM and measuring the perplexity on input $\mathbf{X}$. The highest-scoring $k$ individuals are then chosen as parents for the next evolutionary phase. Nevertheless, the enormous size of the search space means that random application of mutation and crossover can yield suboptimal candidates, resulting in many redundant perplexity evaluations. This inefficiency becomes especially pronounced for large $L'$, as each perplexity computation incurs significant computational overhead.

To mitigate this issue, a monotonicity constraint is enforced such that the rescaled RoPE factors must be non-decreasing: $\lambda_i \leq \lambda_{i+1}$. Only RoPE configurations meeting this criterion are utilized within the LLM during perplexity assessment, thereby minimizing the computational resources required for searching. Concretely, we mandate that $\lambda_i$ follows a strictly increasing trend across the RoPE dimensions ($i=0,...,63$). This monotonicity aligns with the Neural Tangent Kernel (NTK) framework~\cite{jacot2018neural,tancik2020fourier,ntk}, which posits that dimensions associated with higher frequencies (i.e., lower indices) should undergo minimal interpolation (yielding smaller $\lambda_i$), while dimensions tied to lower frequencies (higher indices) permit greater interpolation (resulting in larger $\lambda_i$).

\begin{algorithm}[t]
	\small
	\caption{The search algorithm}
	\textbf{Input:} target LLM, input samples $\mathbf{X}$, population size $P$, mutation size $N_1$, crossover size $N_2$, maximum iterations $\mathcal{T}$, mutation probability $p$ \\

	\begin{algorithmic}[1]
		\label{alg:search}
		\STATE Top-k $\gets \phi$;	
		\STATE $\text{P}_0 \gets \textit{Initialize\_population\_with\_optimization}(P, \mathbf{X}, p)$;
		\FOR{$i$ from $1$ to $\mathcal{T}$}
		\STATE \textit{Compute\_perplexity}(LLM, $\text{P}_{i-1}$, $\mathbf{X}$);
		\STATE Top-k $\gets$ \textit{Update\_Topk}(Top-k, $\text{P}_{i-1}$);
		\STATE $\text{P}_{mutation} \gets \textit{Mutation\_with\_mono\_constraint}(Top-k, N_1, p)$;
		\STATE $\text{P}_{crossover} \gets \textit{Crossover\_with\_mono\_constraint}(Top-k, N_2)$;
		\STATE $\text{P}_i \gets \text{P}_{mutation} \cup \text{P}_{crossover} \cup$ Top-k;
		\ENDFOR
		\STATE Output the member with minimal perplexity from Top-k;
	\end{algorithmic}
\end{algorithm}

\textbf{8$\times$ context window expansion without additional fine-tuning}. Through evolutionary search, we discover non-uniform RoPE rescale factors that strategically maintain important dimensions and token positions, thereby reducing information loss from interpolation. Figure~\ref{fig:non-ft} illustrates that our approach successfully broadens LLaMA2's context range from 4k up to 32k without requiring any fine-tuning. In comparison, other techniques—including PI as well as non-uniform NTK and YaRN—lead to significant increases in perplexity once the window is doubled.

\subsection{Extending LLM Context Window to 2048K}
\label{sec:extendto2048k}

\textbf{Progressive extension to 2048k}. In this section, we present our approach for scaling up the context window of pre-trained LLMs beyond the conventional 4k tokens, reaching up to 2048k. Our proposed non-uniform positional interpolation is capable of achieving an 8$\times$ extension without the need for fine-tuning. For even greater expansions, such as 512$\times$, fine-tuning becomes unavoidable. One possible strategy involves identifying appropriate RoPE rescaling parameters for the desired 2048k context length prior to fine-tuning. Nonetheless, this solution is hindered by the immense computational cost required for training at such scales. Furthermore, our observations indicate that it is notably difficult to effectively fine-tune LLMs under extremely large extension ratios (refer to Appendix).

Fortunately, LongRoPE demonstrates efficacy with both pretrained and further extended LLMs. Accordingly, we propose a streamlined and incremental approach that reaches the 2048k target utilizing only 1k fine-tuning iterations, all conducted at a training sequence length limited to 256k.

$\Diamond$ \textit{Applying LongRoPE search to scale up pre-trained LLMs to 256k}. Using LLaMA2 for illustration, we perform a search procedure aimed at expanding the context window to 128k and 256k. At this point, the context length is increased by factors of 32$\times$ and 64$\times$, respectively.

$\Diamond$ \textit{Fine-tuning to 256k}. Subsequently, we adapt the pre-trained LLM so that it supports a 256k context window. To accomplish this, we initially fine-tune LLaMA2 for 400 iterations with RoPE rescaled factors set for 128k. After completing this stage, we update the RoPE rescaled factors to match 256k on the resulting checkpoint and continue fine-tuning for a further 600 steps. This two-stage approach demonstrates greater efficiency compared to attempting to fine-tune directly to 256k.

$\Diamond$ \textit{Scaling fine-tuned extended LLM to 2048k using LongRoPE search}. In the concluding step, we conduct an additional search process on the LLM that was previously fine-tuned for sequences up to 256k tokens. This approach enables us to achieve a significantly expanded context window of 2048k tokens, eliminating the need for any further fine-tuning. The total achieved extension rate is $512\times$.

\textbf{Recovery of original context window performance}. When scaling up to a substantial 2048k context window, we observe diminished results in the initial window range. This degradation is a documented consequence of positional interpolation~\cite{pi}, which compresses position embeddings in the lower-dimensional region of the base window, thus constraining their expressiveness and impairing the model's efficacy. Specifically, at a 512$\times$ extension factor, positions within the standard 4k context window experience considerable overlap and congestion.

To address this issue, we conduct an additional evolutionary search on the extended LLM, tuning the RoPE rescale factors specifically for shorter context sizes such as 4k and 8k. Because positional interpolation is not as critical at these smaller lengths, we lower the upper bound for the searched $\lambda$. At inference time, the LLM automatically adapts the RoPE rescale factors relevant to the given context length.

\section{LongRoPE}

\label{sec:method}
Building on these insights, we propose LongRoPE, which initially employs an effective search strategy to leverage both identified non-uniformities, and subsequently applies this approach to expand the LLM context window past 2 million tokens.

\subsection{Problem Formulation}

These two forms of non-uniformity result in an expansive solution space and increase the complexity of the optimization process. To mitigate this, we reformulate the multidimensional non-uniform position interpolation optimization as a search problem.

For a LLM targeting a  context window size of $L'$ and lengthy input documents $\mathbf{X}$, where each $\mathbf{x}\in\mathbf{X}$ surpasses $L'$ in token length, we denote the original rotary angle of the $i^{th}$ dimension in RoPE embedding at token position $n$ as  $\frac{n}{\beta^i}$. The optimization problem is then formulated as follows:

\begin{equation}
\begin{gathered}
		\label{eq:problem}

\underset {\mathbf{x}\in\mathbf{X}; \, |\mathbf{x}|\ge L'}{ \operatorname {arg\,min} } \,  \mathcal{L}\, \left( \text{LLM} (\text{RoPE}, \mathbf{X})\right),	\text{where \;}
\\
\scalebox{0.93}{
$\underset {\begin{subarray}{l} i=0,\cdot\cdot,\frac{d}{2}-1;\\ n\in[ 0,|\mathbf{x}|);\end{subarray}}{\text{RoPE}_(n)}= {\left[\cdot\cdot,cos\left(\mathbb{I}(\hat{\lambda}_{i}, \hat{n})\times\frac{n}{\beta^i}\right),  sin\left(\mathbb{I}(\hat{\lambda}_{i}, \hat{n})\times\frac{n}{\beta^i}\right),\cdot\cdot\right] }$}\\
	\text{where \;} \mathbb{I}(\hat{\lambda}_{i},\hat{n})=
	\begin{cases}
	1&\mbox{\; if $n<\hat{n}$}\\
	{\frac{1}{\lambda}_{i}}&\mbox{\; if $n\geq\hat{n}$}
	\end{cases}
	\end{gathered}
\end{equation}

In this framework, we define a collection of scaling coefficients, $\mathbb{I}(\hat{\lambda}_{i},\hat{n})$, which account for both types of non-uniformity. Here, $\hat{\lambda_i}$ represents irregularities specific to RoPE dimensions, while $\hat{n}$ pertains to positional disparities among tokens. More precisely, $\mathbb{I}(\hat{\lambda}_{i},\hat{n})$ is employed to modify the rotational angle for the $i^{th}$ RoPE dimension, utilizing $\hat\lambda_{i}$ as the scaling parameter and $\hat{n}$ as the positional threshold. For the first $\hat{n}$-1 token indices, the adjustment via $\hat\lambda_{i}$ is omitted, and the conventional RoPE rotation angle $\frac{n}{\beta^i}$ remains unchanged. Once token positions reach $n\ge\hat{n}$, the scaling factor is introduced.

Assuming a desired context window length of $L'$, the task is to determine the set of optimal rescale factors—namely, $\mathbb{I}(\hat{\lambda}_{0},\hat{n})$, $\mathbb{I}(\hat{\lambda}_{1},\hat{n})$, ..., $\mathbb{I}(\hat{\lambda}_{i},\hat{n})$, and so forth—across the $1^{st}$ to $d^{th}$ RoPE dimensions. This enables the adapted $\text{LLM}$, incorporating the adjusted $\text{RoPE}$, to achieve the lowest possible next token prediction error, $\mathcal{L}$ (interpreted as perplexity), when processing input sequences $\mathbf{X}$ of length $L'$.

\subsection{Searching the Non-uniform Position Interpolation}

To address the challenge outlined in Eq.~\ref{eq:problem}, we present a straightforward and efficient approach that identifies the optimal RoPE rescale parameters, thereby leveraging the complex and non-uniform characteristics present in multidimensional position embeddings.

\textbf{Search space}. We construct an extensive search space that accommodates the two types of non-uniformity. The search space configuration is summarized in Table~\ref{tbl:searchspace}. In detail, we permit the optimization of a distinct rescale factor for each dimension within the RoPE embedding. To streamline the setup, our approach focuses on tuning $\lambda_i$ and $\hat{n}$, rather than directly exploring $\mathbb{I}(\hat{\lambda}_{i},\hat{n})$, where  $\hat{\lambda}_{i}=1/\lambda_i$. 
According to Table~\ref{tbl:searchspace}, the search for $\lambda_i$ spans from a lower bound of 1.0 (representing pure extrapolation) up to $s\times1.25$ (allowing for interpolation surpassing PI), with increments of 0.01. Here, $s$ denotes the desired expansion ratio of the context window.

The parameter $\hat{n}$ determines how many of the initial token positions maintain the unaltered RoPE embedding, bypassing position interpolation. In practice, we search for optimal values of $\hat{n}$ within the set \{0, 1, 2, 4, 8, 12, 16, 20, 24, 28, 32, 64, 128, 256\}. If $\hat{n}=0$, then every token position employs the learned rescale factors.

\textbf{Evolution-based search}. The search space defined in Table~\ref{tbl:searchspace} includes a vast array of possible positional interpolation configurations, making exhaustive exploration highly impractical. For instance, selecting a $s=4\times$ extension results in $400^{128/2}\times14 = 4\times10^{167}$ potential combinations. As the extension ratio increases, the size of the search space grows exponentially. To cope with this complexity, we employ evolution search~\cite{guo2020single}, and implement two key optimization strategies that significantly enhance the efficiency of the search process. The overall workflow for this approach is detailed in Algorithm~\ref{alg:search}.

\textit{Optimized initial population generation}. Rather than creating an entirely random population of $P$ rescale factors at the outset, our approach incorporates the specific RoPE rescale factors used in PI, NTK, and YaRN directly into the initial set of individuals. The rest of the population, amounting to $P-3$ individuals, is produced by applying random mutations—with probability $p$—to these three established rescale factors.

\textit{Monotonically non-decreasing constraint}. Once the initial population has been produced, we evaluate the LLM perplexity for every individual. This involves applying the designated RoPE rescale factors to the chosen LLM and then determining the perplexity for the input $\mathbf{X}$. The top-k individuals, ranked by their performance, are selected as the parent pool for the evolutionary process. Nevertheless, due to the high dimensionality of the search space, conventional mutation and crossover operations often result in suboptimal candidates, necessitating excessive perplexity evaluations. This challenge becomes even more significant as $L'$ increases, since each perplexity computation is a computationally intensive inference step.

To mitigate this issue, we enforce a non-decreasing monotonicity requirement on the set of RoPE rescaling factors being sampled: $\lambda_i \leq \lambda_{i+1}$. Only those RoPE configurations meeting this monotonicity condition are utilized for perplexity assessment in LLMs, leading to a notable reduction in computational search overhead. In detail, we stipulate that $\lambda_i$ must consistently increase as the RoPE dimension grows, where $i = 0, \ldots, 63$. This monotonic trend across dimensions is motivated by Neural Tangent Kernel (NTK) theory~\cite{jacot2018neural, tancik2020fourier, ntk}, which indicates that lower-dimensional components—associated with higher frequency—necessitate minimal interpolation (hence, a smaller $\lambda_i$), whereas higher-dimensional components—corresponding to lower frequency—permit greater interpolation (implying a larger $\lambda_i$).

\begin{algorithm}[t]
	\small
	\caption{The search algorithm}
	\textbf{Input:} target LLM, input samples $\mathbf{X}$,   population size $P$, mutation size $N_1$, crossover size $N_2$, max iterations $\mathcal{T}$, mutate probability $p$\\

	\begin{algorithmic}[1]
		\label{alg:search}
		\STATE Top-k is initialized as $\phi$;	
		\STATE $\text{P}_0$ is created via \textit{Initialize\_population\_with\_optimization} ($P$, $\mathbf{X}$, $p$); 
		\FOR{$i$ ranging from $1$ to $\mathcal{T}$}
		\STATE \textit{Compute\_perplexity} (LLM, $\text{P}_{i-1}$, $\mathbf{X}$);
		\STATE  Top-k is updated using \textit{Update\_Topk} (Top-k, $\text{P}_{i-1}$);
		\STATE$\text{P}_{mutation}$ is obtained via \textit{Mutation\_with\_mono\_constraint} (Top-k, $N_1$, $p$); 
		\STATE $\text{P}_{crossover}$ is produced by \textit{Crossover\_with\_mono\_constraint} (Top-k, $N_2$);
		\STATE $\text{P}_i$ is defined as the union of $\text{P}_{mutation}$, $\text{P}_{crossover}$, and Top-k;
		\ENDFOR
		\STATE The member in Top-k exhibiting the minimal perplexity is returned;
	\end{algorithmic}
\end{algorithm}

\textbf{8$\times$ extension without fine-tuning}. Leveraging our evolutionary search approach, we discover non-uniform RoPE rescale factors that selectively protect important dimensions and positions, thereby reducing information loss from interpolation. As illustrated in Fig.\ref{fig:non-ft}, this strategy enables us to increase LLaMA2's context window from 4k up to 32k tokens without requiring any fine-tuning. By comparison, prior approaches like PI, and non-uniform NTK and YaRN, exhibit sharp rises in perplexity once the context window is doubled.

\subsection{Extending LLM Context Window to 2048K}
\label{sec:extendto2048k}

\textbf{Progressive extension to 2048k}. Here, we present our approach for scaling the context window of pre-trained LLMs far beyond the typical 4k limit, reaching up to and exceeding 2048k tokens. Utilizing non-uniform positional interpolation enables an 8$\times$ extension with no requirement for fine-tuning. However, when much greater expansion is necessary (such as 512$\times$), fine-tuning becomes unavoidable. One possible strategy involves identifying optimal RoPE rescaling factors tailored to the desired 2048k context length, followed by model fine-tuning. Nevertheless, this is hindered by the exceptionally high computational costs. Furthermore, our empirical results indicate that effectively fine-tuning LLMs at such extensive scaling factors is particularly challenging (refer to Appendix).

Fortunately, LongRoPE proves to be beneficial for both the base model and those LLMs that have undergone extended fine-tuning. As a result, we propose a streamlined, stepwise approach that allows us to reach the desired 2048k target using only 1,000 fine-tuning iterations while keeping the training sequence length constrained to 256k.

$\Diamond$ \textit{Investigating the scaling of pre-trained LLMs to 256k via LongRoPE search}. Using LLaMA2 as the reference model, we perform a search aiming at expanding the context window to 128k and 256k tokens. At this point, the increase in length corresponds to extension factors of 32$\times$ and 64$\times$.

$\Diamond$ \textit{Fine-tuning to 256k}. Subsequently, we adapt the pre-trained LLM to accommodate a 256k context window. Our approach involves initially fine-tuning LLaMA2 for 400 steps utilizing the RoPE rescaling parameters corresponding to 128k. After this stage, we update the RoPE rescaling parameters to those appropriate for 256k at the checkpoint and proceed with a further 600 steps of fine-tuning. This sequential strategy demonstrates greater efficiency compared to performing fine-tuning for 256k in a single stage.

$\Diamond$ \textit{Applying LongRoPE search to expand fine-tuned extended LLM to 2048k.} In the final stage, we conduct an additional search on the LLM that was previously fine-tuned for a 256k sequence length. This approach allows us to attain a substantially larger context window of 2048k, eliminating the need for extra fine-tuning steps. As a result, the extension ratio reaches $512\times$.

\textbf{Restoration of performance in shorter context windows}.  When expanding the context window to a length of 2048k, a decline in performance is observed within the previously standard context window. This drawback is characteristic of positional interpolation methods~\cite{pi}, where embedding positions from the initial context window are compressed into a limited high-dimensional space, ultimately impeding the language model’s effectiveness. In particular, applying an extension ratio of 512$\times$ makes the representation of positions in the original 4k context window highly congested.

To address this issue, an additional evolutionary search is conducted on the expanded LLM to optimize RoPE rescale factors specifically for shorter context lengths such as 4k and 8k. The upper bound for the search over $\lambda$ values is decreased, reflecting the reduced need for positional interpolation at these lengths. The LLM then adapts the relevant RoPE rescale factors on-the-fly during inference.

 \section{Experiments}

\label{sec:exp}
\subsection{Setup}
\textbf{Evaluation Tasks and models}. Our study integrates LongRoPE into both LLaMA2-7B and Mistral-7B models, assessing their performance across three dimensions: (1) the perplexity of language models with extended context windows when processing lengthy documents; (2) the Passkey retrieval task, designed to test the capability of models to identify and extract a straightforward passkey embedded within extensive distractor text; and (3) traditional LLM benchmark evaluations, using a standard context window of 4096 tokens.

\textbf{Fine-tuning}. For LLaMA2, we adopt a linear-decaying learning rate set at 2e-5, paired with a global batch size of 32. Our initial fine-tuning stage comprises 400 steps on the Redpajama~\cite{redpajama} dataset, which we divide into 128k-length chunks marked by BOS and EOS tokens at the boundaries. Subsequently, using the checkpoint from this stage, we conduct 600 additional steps to extend the context window to 256k tokens. Training with the 128k context size utilizes 8 A100 GPUs and employs the distributed system~\cite{cube}, while scaling to 256k tokens requires the use of 16 A100 GPUs. For Mistral, we fix the learning rate at 1e-6 and employ a global batch size of 64 throughout. In both the 128k and 256k configurations, our procedure follows the YaRN~\cite{yarn} methodology: we fine-tune for 400 steps using a sequence length of 16k tokens on Together Computer's Long-Data Collections~\cite{mistraldata}. This process is carried out using 4 A100 GPUs.

\textbf{Search}. For target window sizes up to 256k, we set $P$=64, $N_1$=$N_2$=16, $p$=0.3, and $\mathcal{T}$=40, and in each generation, the top 32 individuals are selected for mutation and crossover. Perplexity assessments use 5 randomly drawn validation examples from the PG19 set, ensuring that each meets or exceeds the desired context length. When targeting windows beyond 512k, we reduce the sizes for population, mutation, and crossover by 50\%. In these cases, perplexity is computed over 3 random validation samples from the Pile-Books3~\cite{pile} dataset.

\textbf{Baselines}. To achieve a context window of 2048k, models were fine-tuned with 128k and 256k window lengths. The resulting configurations, LongRoPE-2048k (ft=128k) for LLaMA2 and LongRoPE-2048k (ft=256k) for Mistral, serve as benchmarks. We evaluate these four models against leading approaches for context window extension, focusing on open-source LLMs refined post-positional interpolation utilizing PI, NTK, and YaRN methods. Among these baselines are Together-32k~\cite{together}, Code LLaMA~\cite{codellama}, LongLoRA-full-FT-100k~\cite{longlora}, as well as YaRN-LLaMA and YaRN-Mistral~\cite{yarn}.

\subsection{Main Results}

\label{sec:mainresult}
\textbf{Language modeling over 256k-token sequences}. Our initial focus is to benchmark performance against leading large language models that support extended context windows, restricting evaluations to a length of 256k tokens. To demonstrate robustness across tasks, we leverage two benchmarks: the Proof-pile~\cite{pg19} and PG19~\cite{pile} test splits. We compute perplexity over various context sizes, employing a moving evaluation window of length 256. On PG19, assessment is conducted on its entire test set consisting of 100 documents. For Proof-pile, following the approach of YaRN~\cite{yarn}, we randomly sample 10 examples, each spanning a minimum of 128k tokens.

Table~\ref{tbl:proof-pile} and Table~\ref{tbl:pg19} present a comparison of the perplexity scores for LLaMA2 and Mistral models, which have been enhanced using various interpolation techniques, evaluated on Proof-pile and PG19 datasets. Two significant findings emerge: \textbf{(1)} The extended models consistently demonstrate reduced perplexity as evaluation length increases from 4k to 256k, indicating improved ability to process and utilize extended contexts. \textbf{(2)} Despite operating within a context window that is 16$\times$ larger—a situation often associated with a drop in performance for shorter sequence lengths—our LongRoPE-2048k models surpass current leading baselines within the 256k context length interval.

\textbf{Long sequence language modeling beyond 2000k}. To assess performance on very lengthy texts, we employ the Books3~\cite{pile} dataset. For practical evaluation, we randomly choose 20 books with lengths greater than 2048k and apply a sliding window of 256k tokens.

Table~\ref{tbl:books3} demonstrates that LongRoPE is effective in extending the context window of both LLaMA2-7B and Mistral-7B models up to 2048k, while maintaining perplexity on par with or better than baseline methods for shorter windows ranging from 8k to 128k. Additionally, the results reveal a clear discrepancy in performance between the 2048k extensions for LLaMA2 and Mistral. Specifically, Mistral surpasses baselines at short context lengths but experiences perplexity greater than 7 for windows larger than 256k. As anticipated, LLaMA2’s perplexity generally declines with increasing context length, though slight increases occur at 1024k and 2048k. Furthermore, for LLaMA2, training LongRoPE-2048k using a 256k fine-tuning length yields better results compared to using 128k, which can be attributed to the lower secondary extension ratio (8$\times$ versus 16$\times$). In contrast, Mistral achieves stronger performance when fine-tuned with a 128k context window. This difference is primarily because, following YaRN's protocol, Mistral’s 128k and 256k fine-tuning leverages a 16k training length, which constrains its capacity to expand the context window after fine-tuning.

\textbf{Passkey retrieval}. Next, we investigate how the size of the effective context window influences generative tasks. Our approach adopts the synthetic passkey retrieval benchmark as introduced by~\cite{passkey}. In this setup, the model must locate and recover a five-digit number, referred to as the passkey, which is embedded at a random point within an extended document. Details of the prompt format are provided in the appendix. The experiment is repeated over 10 trials, each time placing the passkey at a uniformly random position within the total evaluation context length.

Figure~\ref{fig:passkey} presents a comparison of retrieval accuracy across different baseline models. Most prior approaches experience a precipitous decline in accuracy, reaching zero once the token count exceeds 128k. By contrast, LongRoPE-LLaMA2-2048k (ft=256k), even when tasked with extracting a passkey from among millions of tokens, consistently achieves high retrieval accuracy ($\ge$90\%) from 4k up to 2048k. Similarly, LongRoPE-Mistral-2048k (ft=128k) sustains perfect accuracy (100\%) until 1800k tokens and then falls to 60\% at 2048k. This drop is consistent with the observations in Table~\ref{tbl:books3}, where perplexity experiences a modest increase at 2048k.

\textbf{Evaluation on established benchmarks in the native context window}. LongRoPE-2048k models are assessed within their initial context length using the Hugging Face Open LLM Leaderboard~\cite{huggingfaceboard}, applying both zero-shot and few-shot configurations. Specifically, we employ a 25-shot protocol for ARC-Challenge~\cite{allenai:arc}, a 10-shot setup for HellaSwag~\cite{zellers2019hellaswag}, a 5-shot format for MMLU~\cite{mmlu}, and a zero-shot approach for TruthfulQA~\cite{lin2021truthfulqa}.

As indicated in Table~\ref{tbl:commonsense}, our models attain similar performance levels on the standard benchmark, which was originally intended for shorter context lengths, and surpass the original Mistral model by +0.5\% on the TruthfulQA evaluation. Although LongRoPE-LLaMA2-2048k, fine-tuned at 256k, exhibits a marginally greater decline in performance, its results across most tasks remain within acceptable limits.

\subsection{Ablation Results}

\textbf{Effectiveness of the second positional interpolation}. Within our progressive extension framework, we employ our search algorithm to perform a second application of non-uniform positional interpolation on the extended, fine-tuned LLMs. We assess this approach using experiments conducted on our fine-tuned LLaMA2-256k model, which we further scale to 512k, 1024k, and 2048k using PI and YaRN methods. As demonstrated in Table~\ref{tbl:secondsearch}, our method of non-uniform positional interpolation maintains stable perplexity across these extension scales. In comparison, both PI and YaRN show rapid increases in perplexity as the extension ratio grows.

\textbf{Recovery performance at reduced context lengths.} In order to address potential declines in accuracy when handling shorter context windows, we re-optimized the RoPE scaling parameters for LongRoPE-2048k by applying our search methodology. This entailed constraining the maximum scale values available during the search process, which discourages excessive interpolation specifically at the 4k and 8k context window sizes. In Table~\ref{tbl:shortperformance}, we present a perplexity evaluation of LongRoPE-LLaMA2-2048k on the Proof-pile dataset for 4k and 8k contexts, supplemented by the average LLM benchmark accuracy. These findings provide clear evidence of notable gains in performance at shorter context lengths.

\textbf{Analysis of the two types of non-uniformity}. To assess the influence of each non-uniformity on model performance, we conduct an ablation study. Two sets of experiments are devised: (i) scaling LLaMA2-7B to shorter contexts (16k and 32k) using three strategies—PI, optimizing the RoPE dimension alone, and optimizing both forms of non-uniformity; (ii) expanding our fine-tuned LLaMA2 model from 256k to 2048k sequence length, following the analogous approach. Perplexity measurements are obtained prior to any fine-tuning. Table~\ref{tbl:searchdimension} illustrates that introducing non-uniformity in the RoPE dimension yields substantial reductions in perplexity when compared with PI's linear interpolation. While non-uniformity in token position provides notable gains at sequence lengths of 16k and 32k, this benefit appears diminished at 2048k, potentially due to the excessive sequence length. Simply retaining original tokens without interpolation proves ineffective in this context; further investigation is deferred to future work.

\section{Related Works}

Beyond position interpolation techniques, this section reviews other categories of related literature.

\textbf{Retrieval-based} methods incorporate external memory components to store extended historical context and utilize retrieval mechanisms to access relevant documents during inference~\cite{longllama,wang2023augmenting,borgeaud2022improving}. Such strategies frequently require direct architectural changes to the LLM models. In contrast, our approach remains relatively lightweight, involving only minor alterations to positional embeddings. Moreover, our method supports a broader range of long-context applications extending beyond retrieval scenarios, including tasks like long document summarization and few-shot learning.

\textbf{Attention-based context window extensions}. In addition to approaches like positional embedding interpolation, other studies have managed to extend input context lengths while retaining the original LLM context window size by altering the underlying attention mechanisms~\cite{han2023lminfinite, streamingllm,parallelwindow}. Central to these techniques is the introduction of innovative attention masks that address the attention explosion problem introduced by novel position indices. Notably, these attention-based strategies can be used alongside positional interpolation methods, as they provide complementary benefits.

\textbf{Fine-tuning based approaches} investigate strategies for adapting pre-trained LLMs to handle extended contexts by altering position embeddings. Prior research, including Code LLaMA~\cite{codellama}, LLaMA2 Long~\cite{xiong2023effective}, and ScaledRoPE~\cite{liu2023scaling}, typically employ an enlarged base value for RoPE and perform fine-tuning at the intended sequence length. In contrast, our approach is adaptable to a range of target lengths and is capable of supporting contexts exceeding 2M tokens. As extending fine-tuning to long contexts (over 128k tokens) significantly increases GPU consumption, methods such as LongLoRA~\cite{longlora} and PoSE~\cite{zhu2023pose} have been developed to address computational costs. Our technique is complementary and independent of these efficient fine-tuning solutions.

\section{Conclusion}

In this paper, we introduce LongRoPE, a technique that significantly increases the context window of LLMs to an extraordinary 2048k, all while preserving their performance on the original, shorter contexts. Our approach leverages two distinct types of non-uniformities within RoPE positional embeddings through a streamlined evolutionary search process. This methodology yields dual advantages: it results in strong initialization for subsequent fine-tuning and facilitates an 8$\times$ enlargement of the context window even in the absence of fine-tuning. Furthermore, we implement a stepwise extension scheme that employs LLMs fine-tuned at 256k context length to attain a 2048k context window, eliminating the need for additional fine-tuning stages. Thorough experimentation demonstrates LongRoPE’s efficacy. We anticipate that LongRoPE-2048k will support numerous novel applications requiring extended context lengths and encourage continued research in this area.

\section*{Broader Impacts} The objective of this research is to contribute to the progress of Machine Learning as a discipline. While various societal effects may arise from our findings, we do not identify any particular consequence that warrants special attention in this context.

	\balance

\end{document}